\sgasection{6.4}{Products of root data}
\begin{definition}{6.4.1}\label{XXI.6.4.1}
Let $\mathcal R=(M,M^*,R,R^*)$ and $\mathcal R'=(M',M'^*,R',R'^*)$ be two root data. One calls the \emph{product root datum} of $\mathcal R$ and $\mathcal R'$, and denotes by $\mathcal R''=\mathcal R\times\mathcal R'$, the root datum $(M'',M''^*,R'',R''^*)$ where
\[
\begin{gathered}
M''=M\times M',\qquad M''^*=M^*\times M'^*,\\
R''=(R\times0)\cup(0\times R'),\qquad R''^*=(R^*\times0)\cup(0\times R'^*),
\end{gathered}
\]
the map $\alpha\mto\alpha^*$ being the obvious map.
\end{definition}
\begin{proposition}{6.4.2}\label{XXI.6.4.2}
Under the preceding conditions, one has canonical isomorphisms
\[
\Gamma_0(R'')\simeq\Gamma_0(R)\times\Gamma_0(R'),\qquad
\mathcal V(R'')\simeq\mathcal V(R)\times\mathcal V(R'),
\]
\[
W(\mathcal R'')\simeq W(\mathcal R)\times W(\mathcal R'),
\]
etc., and the equalities
\[
\begin{gathered}
\operatorname{rgred}(\mathcal R'')=\operatorname{rgred}(\mathcal R)+\operatorname{rgred}(\mathcal R'),\\
\operatorname{rgss}(\mathcal R'')=\operatorname{rgss}(\mathcal R)+\operatorname{rgss}(\mathcal R').
\end{gathered}
\]
One also has a canonical isomorphism of root data
\[
(\mathcal R\times\mathcal R')^*\simeq\mathcal R^*\times\mathcal R'^*.
\]
\end{proposition}
The preceding definitions generalize at once to a product of several factors. One has at once:
\begin{proposition}{6.4.3}\label{XXI.6.4.3}
Let $\mathcal R=\mathcal R_1\times\mathcal R_2\times\cdots\times\mathcal R_n$ be a product of root data. The following conditions are equivalent:
\begin{enumeratei}
\item $\mathcal R$ is semisimple (resp. simply connected, resp. adjoint, resp. reduced).
\item Each $\mathcal R_i$ is semisimple (resp. simply connected, resp. adjoint, resp. reduced).
\end{enumeratei}
\end{proposition}
Consider the following particular case: let $\mathcal R_0$ be a trivial root datum and $\mathcal R_1$ a semisimple root datum. One then has a commutative diagram
\[
\xymatrix{&\mathcal R_0\times\mathcal R_1\ar[dr]&\\
\mathcal R_1\ar[ur]\ar[rr]_{\mathrm{id}}&&\mathcal R_1.}
\]
\begin{lemma}{6.4.4}\label{XXI.6.4.4}
One has canonical isomorphisms
\[
\xymatrix{\operatorname{corad}(\mathcal R_0\times\mathcal R_1)\ar[rr]^{\sim}\ar[dr]_{\simeq}&&\operatorname{rad}(\mathcal R_0\times\mathcal R_1)\\
&\mathcal R_0\ar[ur]_{\simeq}&}.
\]
In particular, $N(\mathcal R_0\times\mathcal R_1)=0$.
\end{lemma}
We shall see later that, conversely, if $N(\mathcal R)=0$, the root datum $\mathcal R$ is the product of a semisimple datum and a trivial datum.

\sgasection{6.5}{Induced and coinduced root data}
Let $\mathcal R=(M,M^*,R,R^*)$ be a root datum. Let $N\subset M$ be a subgroup containing the roots, i.e. such that
\[
\Gamma_0(R)\subset N\subset M.
\]
The canonical linear map $i_N:N\to M$ gives by transposition a linear map
\[
{}^ti_N:M^*\to N^*.
\]
Put $R_N=R$ and $R_N^*={}^ti_N(R^*)$.
\begin{lemma}{6.5.1}\label{XXI.6.5.1}
$\mathcal R_N=(N,N^*,R_N,R_N^*)$ is a root datum, and $i_N$ a morphism.
\end{lemma}
Let us first show that ${}^ti_N$ induces an isomorphism from $R^*$ onto $R_N^*$. If $\alpha,\beta\in R$ and ${}^ti_N(\alpha^*)={}^ti_N(\beta^*)$, one has $(\alpha^*,x)=(\beta^*,x)$ for every $x\in N$, in particular for $x\in R$, which gives $\alpha=\beta$ by 1.1.4. The rest follows without difficulty.
\begin{definition}{6.5.2}\label{XXI.6.5.2}
$\mathcal R_N$ is called the root datum \emph{induced} by $\mathcal R$ on $N$.
\end{definition}
\begin{lemma}{6.5.3}\label{XXI.6.5.3}
Let $f:\mathcal R'\to\mathcal R$ be a morphism. Put $N=f(M')\subset M$. Then $f$ factors uniquely through $i_N$.
\end{lemma}
In particular, isogenies $\mathcal R'\to\mathcal R$, up to isomorphism, correspond bijectively to subgroups of finite index of $M/\Gamma_0(R)$, which makes 6.2.7 more precise.

Let now $N^*$ be a subgroup of $M^*$ containing $R^*$. One defines the datum \emph{coinduced} by $\mathcal R$ on $N^*$ by
\[
\mathcal R^{N^*}=(\mathcal R^*_{N^*})^*,
\]
and one has a canonical morphism:
\[
p^{N^*}:\mathcal R\to\mathcal R^{N^*}.
\]
\begin{lemma}{6.5.4}\label{XXI.6.5.4}
Let $f:\mathcal R'\to\mathcal R$ be a morphism. There exist subgroups $N\subset M$ and $N'^*\subset M'^*$ such that $f$ factors as
\[
\xymatrix{\mathcal R'\ar[r]^f\ar[d]_{p^{N'}}&\mathcal R\\
\mathcal R'^{N'^*}\ar[r]^{f_0}_{\sim}&\mathcal R_N\ar[u]_{i_N},}
\]
where $f_0$ is an isomorphism.
\end{lemma}
Indeed, take $N=f(M')$ as in 6.5.3. The resulting morphism $M'\to N$ is surjective, hence its transpose injective. Take the image of the latter as $N'^*$.

Let us now treat certain particular cases. If one takes $N=\Gamma_0(R)$, one will write $\mathcal R_N=\operatorname{ad}(\mathcal R)$. If one takes $N=\mathcal V(R)\cap M$, one will write $\mathcal R_N=\operatorname{ss}(\mathcal R)$. One therefore has a diagram:
\[
\operatorname{ad}(\mathcal R)\to\operatorname{ss}(\mathcal R)\to\mathcal R.
\]
Put $\operatorname{d\acute er}(\mathcal R)=\operatorname{ss}(\mathcal R^*)^*$ and $\operatorname{sc}(\mathcal R)=\operatorname{ad}(\mathcal R^*)^*$; by duality, one obtains a diagram:
\[
\mathcal R\to\operatorname{d\acute er}(\mathcal R)\to\operatorname{sc}(\mathcal R).
\]
\begin{proposition}{6.5.5}\label{XXI.6.5.5}
\begin{enumeratei}
\item In the first row of the diagram
\[
\xymatrix@C=1.5em{
\operatorname{ad}(\mathcal R)\ar[r]&\operatorname{ss}(\mathcal R)\ar[rr]\ar[dr]&&\operatorname{d\acute er}(\mathcal R)\ar[r]&\operatorname{sc}(\mathcal R)\\
&&\mathcal R\ar[ur]&&
}
\]
the four data are semisimple and the three morphisms are isogenies.
\item $\operatorname{ad}(\mathcal R)$ is an adjoint datum, and $\mathcal R$ is adjoint if and only if $\operatorname{ad}(\mathcal R)\to\mathcal R$ is an isomorphism.
\item $\operatorname{sc}(\mathcal R)$ is a simply connected datum, and $\mathcal R$ is simply connected if and only if $\mathcal R\to\operatorname{sc}(\mathcal R)$ is an isomorphism.
\item The following conditions are equivalent:
\begin{enumeratea}
\item $\mathcal R$ is semisimple,
\item $\operatorname{ss}(\mathcal R)\to\mathcal R$ is an isomorphism,
\item $\mathcal R\to\operatorname{d\acute er}(\mathcal R)$ is an isomorphism.
\end{enumeratea}
\end{enumeratei}
\end{proposition}
Let us pause a moment on the morphism $\operatorname{ss}(\mathcal R)\to\operatorname{d\acute er}(\mathcal R)$. Returning to the construction of $\operatorname{ss}(\mathcal R)$ and $\operatorname{d\acute er}(\mathcal R)$, it is easy to prove:
\begin{lemma}{6.5.6}\label{XXI.6.5.6}
Let $h:\operatorname{ss}(\mathcal R)\to\operatorname{d\acute er}(\mathcal R)$ be the canonical isogeny. One has $K(h)\simeq N(\mathcal R)$.
\end{lemma}

\numpara{6.5.7} Consider other particular cases of induced data. Put
\[
N=\{x\in M\mid(\alpha^*,x)=0\text{ for }\alpha\in R\}\times\Gamma_0(R);
\]
one knows that the sum is direct by 1.2.5. It follows that the root datum $\mathcal R_N$ is identified with the product $\operatorname{ad}(\mathcal R)\times\operatorname{corad}(\mathcal R)$.

One can do the same replacing $\Gamma_0(R)$ by $\mathcal V(R)\cap M$, then dualize these two constructions. One thus obtains a diagram of root data:
\[
\resizebox{0.97\linewidth}{!}{$\xymatrix@C=0.8em{
\operatorname{ad}(\mathcal R)\ar[r]\ar[d]&\operatorname{ss}(\mathcal R)\ar[rr]\ar[d]&&\operatorname{d\acute er}(\mathcal R)\ar[r]\ar[d]&\operatorname{sc}(\mathcal R)\ar[d]\\
\operatorname{ad}(\mathcal R)\times\operatorname{corad}(\mathcal R)\ar[r]\ar[d]&\operatorname{ss}(\mathcal R)\times\operatorname{corad}(\mathcal R)\ar[r]\ar[d]&\mathcal R\ar[r]&\operatorname{d\acute er}(\mathcal R)\times\operatorname{rad}(\mathcal R)\ar[r]\ar[d]&\operatorname{sc}(\mathcal R)\times\operatorname{rad}(\mathcal R)\ar[d]\\
\operatorname{ad}(\mathcal R)\ar[r]&\operatorname{ss}(\mathcal R)\ar[rr]&&\operatorname{d\acute er}(\mathcal R)\ar[r]&\operatorname{sc}(\mathcal R)
}$};
\]
which is commutative, as one checks at once. This diagram is self-dual in an obvious sense. The horizontal morphisms are isogenies. The composites of the vertical arrows are the identity.
\begin{lemma}{6.5.8}\label{XXI.6.5.8}
Let $h_1$ and $h_2$ be the canonical isogenies:
\[
\operatorname{ss}(\mathcal R)\times\operatorname{corad}(\mathcal R)\xrightarrow{h_1}\mathcal R\xrightarrow{h_2}\operatorname{d\acute er}(\mathcal R)\times\operatorname{rad}(\mathcal R).
\]
One has $K(h_1)\simeq K(h_2)\simeq N(\mathcal R)$.
\end{lemma}
This is trivial from the definitions.
\begin{corollary}{6.5.9}\label{XXI.6.5.9}
Let $\mathcal R$ be a root datum. The following conditions are equivalent:
\begin{enumeratei}
\item $N(\mathcal R)=0$, i.e. $\operatorname{corad}(\mathcal R)\to\operatorname{rad}(\mathcal R)$ is an isomorphism.
\item $h:\operatorname{ss}(\mathcal R)\to\operatorname{d\acute er}(\mathcal R)$ is an isomorphism.
\item $h_1:\operatorname{ss}(\mathcal R)\times\operatorname{corad}(\mathcal R)\to\mathcal R$ is an isomorphism.
\item $h_2:\mathcal R\to\operatorname{d\acute er}(\mathcal R)\times\operatorname{rad}(\mathcal R)$ is an isomorphism.
\item $\mathcal R$ is the product of a semisimple datum and a trivial datum.
\end{enumeratei}
\end{corollary}
Let us also state a trivial consequence of the preceding remarks:
\begin{corollary}{6.5.10}\label{XXI.6.5.10}
For every root datum $\mathcal R$, there exist isogenies
\[
\operatorname{ad}(\mathcal R)\times\mathcal R_0\to\mathcal R\to\operatorname{sc}(\mathcal R)\times\mathcal R_0,
\]
where $\mathcal R_0$ is ``the'' trivial root datum of rank $\operatorname{rgred}(\mathcal R)-\operatorname{rgss}(\mathcal R)$.
\end{corollary}
Finally, let us point out a result which may be useful:
\begin{lemma}{6.5.11}\label{XXI.6.5.11}
Let $\mathcal R$ be a root datum, $\Delta$ a system of simple roots, $\Delta'$ a subset of $\Delta$; consider the root datum (cf. 3.4.7)
\[
\mathcal R_{\Delta'}=(M,M^*,R_{\Delta'},R^*_{\Delta'}).
\]
\begin{enumeratei}
\item If $\mathcal R$ is simply connected, $\operatorname{d\acute er}(\mathcal R_{\Delta'})$ is simply connected.
\item If $\mathcal R$ is adjoint, $\operatorname{ss}(\mathcal R_{\Delta'})$ is adjoint.
\end{enumeratei}
\end{lemma}
The two assertions are clearly equivalent by duality. The second reduces to checking the formula:
\[
M\cap\mathcal V(R_{\Delta'})=\Gamma_0(R_{\Delta'});
\]
now, if $M=\Gamma_0(R)$, both sides are equal to the subgroup of $M$ generated by $\Delta'$.

\sgasection{6.6}{Weights}
\begin{definition}{6.6.1}\label{XXI.6.6.1}
Let $\mathcal R$ be a root datum. Put{\renewcommand{\thefootnote}{(26)}\nde{$\Gamma(\mathcal R)$ has been replaced by $\Lambda(\mathcal R)$ to avoid any risk of confusion with $\Gamma_0(R)$.}}
\[
\Lambda(\mathcal R)=\{x\in V(R)\mid(\alpha^*,x)\in\ZZ\text{ for every }\alpha^*\in R^*\}.
\]
The elements of $\Lambda(\mathcal R)$ are called the \emph{weights} of $\mathcal R$. The weights of $\mathcal R^*$ are called the \emph{coweights} of $\mathcal R$.
\end{definition}
One has $\Gamma_0(R)\subset\Lambda(\mathcal R)$, and $\Lambda(\mathcal R)$ is stable under $W(\mathcal R)$.
\begin{lemma}{6.6.2}\label{XXI.6.6.2}
The bilinear map $V^*\times V\to\QQ$ induces a duality
\[
\Gamma_0(R^*)\times\Lambda(\mathcal R)\to\ZZ.
\]
\end{lemma}
Trivial.
\begin{corollary}{6.6.3}\label{XXI.6.6.3}
Let $\Delta^*=(\alpha_1^*,\ldots,\alpha_n^*)$ be a system of simple coroots. Let $p_i$, $i=1,2,\ldots,n$, be the elements of $\mathcal V(R)$ defined by
\[
(\alpha_i^*,p_j)=\delta_{ij},
\]
(whence $s_{\alpha_i}(p_i)=p_i-\alpha_i$ and $s_{\alpha_i}(p_j)=p_j$ for $i\ne j$).{\renewcommand{\thefootnote}{($*$)}\footnote{Attention: if the datum is not reduced, the $\alpha_i$ do not form a system of simple roots.}} Then $\Lambda(\mathcal R)$ is the free abelian group generated by the $p_i$.
\end{corollary}
The $p_i$ are called the \emph{fundamental weights} corresponding to the system of simple coroots $\Delta^*$.
\begin{corollary}{6.6.4}\label{XXI.6.6.4}
For every $\alpha^*\in\Delta^*$, one therefore has $(\alpha^*,\sum_i p_i)=1$, hence $\sum_i p_i=\rho_{R_+}$ (cf. 3.5.1), where $R_+=\mathcal P(\operatorname{ind}(\Delta))$.
\end{corollary}
\begin{corollary}{6.6.5}\label{XXI.6.6.5}
For every $x\in\mathcal V(R)$, one has $x=\sum_i(\alpha_i^*,x)p_i$.
\end{corollary}
Note that $R^*\subset\Gamma_0(R^*)$ and $R\subset\Lambda(\mathcal R)$, hence $(\Lambda(\mathcal R),\Gamma_0(R^*),R,R^*)$ is a root datum.
\begin{corollary}{6.6.6}\label{XXI.6.6.6}
The canonical morphism $\Gamma_0(R^*)\to M^*$ is the transpose of the morphism $x\mto\sum_i(\alpha_i^*,x)p_i$, which defines a morphism of root data, and one has a commutative diagram:
\[
\xymatrix{&\operatorname{sc}(\mathcal R)\ar[d]^{\simeq}\\
\mathcal R\ar[ur]\ar[r]&(\Lambda(\mathcal R),\Gamma_0(R^*),R,R^*).}
\]
\end{corollary}
One therefore has an explicit description of $\operatorname{sc}(\mathcal R)$ in terms of the weights of $\mathcal R$. Similarly, one finds a commutative diagram:
\[
\xymatrix{\operatorname{ad}(\mathcal R)\ar[d]_{\simeq}\ar[dr]&\\
(\Gamma_0(R),\Lambda(\mathcal R^*),R,R^*)\ar[r]&\mathcal R}.
\]
\begin{corollary}{6.6.7}\label{XXI.6.6.7}
For $\mathcal R$ to be simply connected, it is necessary and sufficient that $M=\Lambda(\mathcal R)$.
\end{corollary}
\begin{remark}{6.6.8}\label{XXI.6.6.8}
One has $\Lambda(\mathcal R)\cap M=\mathcal V(R)\cap M$. For $\mathcal R$ to be semisimple, it is therefore necessary and sufficient that $M\subset\Lambda(\mathcal R)$.
\end{remark}
It also follows from the results of 6.5 that:
\begin{corollary}{6.6.9}\label{XXI.6.6.9}
For $\mathcal R$ to be the product of a simply connected datum and a trivial datum, it is necessary and sufficient that $M\supset\Lambda(\mathcal R)$.
\end{corollary}
Consider now the canonical isogeny
\[
f:\operatorname{ad}(\mathcal R)\to\operatorname{sc}(\mathcal R),
\]
and put $Z(\mathcal R)=K(f)$. One has $Z(\mathcal R)\simeq Z(\operatorname{sc}\mathcal R)\simeq Z(\operatorname{ad}\mathcal R)$.
\begin{corollary}{6.6.10}\label{XXI.6.6.10}
One has a canonical isomorphism $Z(\mathcal R)=\Lambda(\mathcal R)/\Gamma_0(R)$. More precisely, one has an exact sequence of $W(\mathcal R)$-modules:
\[
0\to\Gamma_0(R)\to\Lambda(\mathcal R)\to Z(\mathcal R)\to0.
\]
\end{corollary}
\begin{corollary}{6.6.11}\label{XXI.6.6.11}
One has a canonical pairing
\[
Z(\mathcal R^*)\times Z(\mathcal R)\to\QQ/\ZZ
\]
putting these groups in duality.
\end{corollary}
\begin{remark}{6.6.12}\label{XXI.6.6.12}
One has $Z(\mathcal R\times\mathcal R')\simeq Z(\mathcal R)\times Z(\mathcal R')$. Consider, in particular, simply connected data $\mathcal R_i$, $i=1,2,\ldots,n$, and a trivial datum $\mathcal R_0$. Put $\mathcal R=\mathcal R_0\times\mathcal R_1\times\cdots\times\mathcal R_n$. Let $\mathcal R=(M,M^*,R,R^*)$, $\mathcal R_0=(M_0^*,M_0^*,\varnothing,\varnothing)$. One has
% typo?: repeated M_0 star in the trivial datum kept as printed.
\[
M/\Gamma_0(R)\simeq M_0\times Z(\mathcal R_1)\times\cdots\times Z(\mathcal R_n).
\]
\end{remark}

\sgasection{6.7}{Automorphisms}
An automorphism of $\mathcal R$ is, by 6.1.4, an automorphism $u$ of $M$ such that $u(R)=R$, ${}^tu(R^*)=R^*$. In particular, every element $w$ of $W(\mathcal R)$ defines an automorphism of $\mathcal R$.
\begin{lemma}{6.7.1}\label{XXI.6.7.1}
$W(\mathcal R)$ is an invariant subgroup of $\Aut(\mathcal R)$. More precisely, if $u\in\Aut(\mathcal R)$ and $\alpha\in R$, one has
\[
us_\alpha u^{-1}=s_{u(\alpha)}.
\]
\end{lemma}
